\begin{trapbox}
	\textbf{\textcolor{CTrap}{易错点一：只看公式，不看定义域}}

	判断奇偶性要先确认：若 $x$ 在函数定义域内，
	$-x$ 也必须在定义域内。若原函数已知为奇函数或偶函数，
	其定义域本身就满足这一条件；做除法时，还要重新确认
	商函数的定义域。

	\medskip
	\textbf{\textcolor{CTrap}{易错点二：认为分母必须处处不为零}}

	商函数只在分母不为零的点上定义。
	例如 $x/x$ 的定义域为 $x\ne0$，虽然分母在 $0$ 处为零，
	商函数仍然可以判断奇偶性。关键是删去零点后的定义域
	关于原点对称。

	\medskip
	\textbf{\textcolor{CTrap}{易错点三：把奇乘奇记成奇}}

	计算 $f(-x)g(-x)$：若两函数都为奇函数，
	就会出现两个负号，
	\[
		(-f(x))(-g(x))=f(x)g(x).
	\]
	结果不变，所以奇乘奇为偶。除法同理：奇除奇为偶。

	\medskip
	\textbf{\textcolor{CTrap}{易错点四：把奇加偶判成固定类型}}

	例如 $x+x^2$ 通常既非奇函数也非偶函数。
	但不能一概说“奇加偶一定非奇非偶”：
	零函数既是奇函数也是偶函数，若其中一项为零，
	结果可能仍为奇函数或偶函数。遇到这种情形，
	直接计算 $F(-x)$ 判断。
\end{trapbox}